% ---------------------------------------------------------------------------
% Front matter: title, authors, graphical abstract, abstract, key points.
%
% Framing follows the project objective in CLAUDE.md: this is a FEASIBILITY STUDY.
% The question is whether DNAGPT inference can run under homomorphic encryption, and where
% correctness, performance, or memory would stop a complete encrypted deployment.
% Feasibility and practicality are reported as SEPARATE verdicts. A correct-but-slow encrypted
% path is a valid, publishable outcome; nothing is tuned to force "practical".
%
% Source: context/01_scenario_and_motivation.md, context/05_claims_and_qualifiers.md
% ---------------------------------------------------------------------------

\title{\vspace{-1.2cm}\textbf{Feasibility of Homomorphic Inference for a Genomic Foundation
Model: A Client-Assisted CKKS Evaluation of DNAGPT}\vspace{-0.2cm}}

\author{%
  \begin{minipage}{0.96\textwidth}
  \raggedright
  \textbf{Christos Galanopoulos$^{1,*}$, Kimon Antonios Provatas$^{1}$,
  Ilias Georgakopoulos-Soares$^{1,*}$}\\[2pt]
  \normalsize $^{1}$The University of Texas at Austin, Austin, TX, USA\\[2pt]
  \small $^{*}$Corresponding authors: \texttt{TODO@utexas.edu}\\[2pt]
  \small\textit{Homomorphic encryption $\cdot$ Genomic privacy $\cdot$ Transformer inference
  $\cdot$ GPU acceleration}
  \end{minipage}
}
\date{}

\maketitle
\thispagestyle{fancy}

% The graphical abstract is temporarily excluded pending evidence-safe regeneration; see
% paper-docs/TODO.md. The generated asset contains timing and client-memory claims whose
% supporting logs are not present in the repository.

% --- Abstract ---------------------------------------------------------------
\section*{Abstract}

\textbf{Motivation.} Human genomic sequence is identifying, cannot be replaced after disclosure,
and is precisely the input that genomic foundation models are designed to interpret. We ask
whether a provider can execute a released genomic transformer without receiving genomic values in
plaintext, and whether correctness, memory, or cost prevents a complete encrypted deployment.

\textbf{Results.} We validate DNAGPT on genomic-signal recognition, promoter and splice-site
classification, and mRNA abundance regression, then freeze an independently checked numerical
reference. We evaluate two CKKS designs. In the evaluated pure non-interactive configuration,
polynomial nonlinearities permit a released-weight block to complete, but context and evaluation
keys exceed the tested $80$~GB memory envelope when configured for composition. In the
client-assisted design, the provider evaluates linear operations on ciphertexts while the
key-holding data owner evaluates exact LayerNorm, causal softmax, and GELU at fixed boundaries.
The resulting released-weight block at the task-defined prompt matches its validated plaintext
reference with $4.64\times10^{-9}$ relative error against a $4\times10^{-2}$ tolerance fixed in
advance, while the evaluating process peaks at $9.6$~GiB of GPU memory.

\textbf{Conclusions.} Correctness is not the barrier. The protocol moves the accelerator-memory
barrier by replacing deep non-interactive segments with fresh encryption at declared boundaries;
bounded caching separately prevents encoded plaintexts from accumulating in host memory. The
systems campaign removes repeated host encoding while preserving an asserted operation schedule,
packing, depth, and reference. Client-assisted DNAGPT evaluation is feasible at the demonstrated
scope. Its practicality remains unresolved until the retained implementation is timed on a clean,
dedicated node.

\textbf{Availability.} Code, evaluation harnesses, and the evidence record are maintained at
\url{https://github.com/Georgakopoulos-Soares-lab/private_genomic_dnagpt}.

% --- Key points -------------------------------------------------------------
\section*{Key Points}

\begin{itemize}
  \item We test whether a released genomic foundation model can be evaluated under homomorphic
        encryption, and report feasibility and practicality as separate verdicts rather than as
        one.
  \item The model is validated first on three genomic task families; the independent numerical
        reference is itself asserted against every upstream block before it is used to verify
        encrypted execution.
  \item Correctness is not a barrier: a complete transformer block at the task's own prompt length
        matches the validated plaintext reference to $4.64\times10^{-9}$ relative error, seven
        orders of magnitude inside the tolerance fixed before the run.
  \item Memory is a barrier that design moves. Approximating the nonlinearities under encryption
        exceeded the tested $80$~GB accelerator envelope at composition; client assistance yields
        a $9.6$~GiB evaluating-process GPU peak for the retained block.
  \item Bounded caching addresses a distinct host-memory problem. Runtime assertions confirm that
        the complete operation schedule, packing, depth, and reference remain unchanged while
        repeated plaintext encoding is removed.
  \item The data owner requires no GPU. Practicality remains a separate, unresolved verdict until
        the retained implementation has a clean dedicated-node timing measurement.
\end{itemize}
